\documentclass[11pt,a4paper]{article}

\usepackage[margin=1in]{geometry}
\usepackage{amsmath,amssymb}
\usepackage{booktabs}
\usepackage{tabularx}
\usepackage{array}
\usepackage{enumitem}
\usepackage{hyperref}
\usepackage{microtype}
\usepackage{xcolor}
\usepackage{listings}
\usepackage{caption}

\hypersetup{
  colorlinks=true,
  linkcolor=blue!50!black,
  citecolor=blue!50!black,
  urlcolor=blue!50!black
}

\newcolumntype{Y}{>{\raggedright\arraybackslash}X}

\newcommand{\layer}[1]{\paragraph{#1.}}
\newcommand{\setup}{\layer{Setup}}
\newcommand{\problem}{\layer{Problem}}
\newcommand{\rootcause}{\layer{Root cause}}
\newcommand{\solution}{\layer{Solution}}
\newcommand{\payoff}{\layer{Payoff}}

\title{Large Language Models End to End\\
\large Explained the Interview Way:\\
Basics $\rightarrow$ Problem $\rightarrow$ Root Cause $\rightarrow$ Solution}
\author{Validated conceptual survey\\[0.5em]
Style example:
\href{https://www.youtube.com/watch?v=-AB6m0Spo6c}{Vizuara --- LLM Interview Series \#5: What Is PagedAttention?}}
\date{}

\begin{document}
\maketitle

\begin{abstract}
This document surveys the LLM lifecycle---pretraining, post-training, inference, and evaluation---using an interview-style narration:
do not jump to definitions; first establish basics, then the failure of the naive approach, then the root design cause, and only then the method.
The pedagogical pattern follows the public explanation style in
\href{https://www.youtube.com/watch?v=-AB6m0Spo6c}{Vizuara's PagedAttention interview video},
applied uniformly to every major concept.
Technical claims are anchored to landmark papers and primary writeups; numeric speedups are not treated as universal laws.
\end{abstract}

\tableofcontents
\clearpage

% =============================================================================
\section{The Explanation Contract}
\label{sec:contract}

In the PagedAttention interview video, the narrator states explicitly that answering ``What is PagedAttention?'' by starting with the definition is the wrong move.
A strong answer has layers~\cite{vizuara2026paged}:

\begin{enumerate}[leftmargin=*]
  \item \textbf{Basics / setup} --- what must the listener already understand?
  \item \textbf{Problem} --- what breaks, with a concrete example?
  \item \textbf{Root cause} --- which design choice creates the failure?
  \item \textbf{Solution} --- only then name the method and walk the mechanism.
  \item \textbf{Payoff} --- what improves in a real system, and what tradeoff appears?
\end{enumerate}

\begin{table}[ht]
\centering
\small
\begin{tabularx}{\textwidth}{@{}lY@{}}
\toprule
\textbf{Layer} & \textbf{Question answered} \\
\midrule
Setup & What world are we in? \\
Problem & What fails if we do the naive thing? \\
Root cause & Which assumption creates that failure? \\
Solution & What method addresses that root cause? \\
Payoff & Why practitioners use it \\
\bottomrule
\end{tabularx}
\caption{Narration contract used for every concept in this survey.}
\label{tab:contract}
\end{table}

If you only remember the name of a method, you have not explained it.

% =============================================================================
\section{Foundations}
\label{sec:foundations}

\subsection{What is an LLM, really?}
\label{sec:what-llm}

\setup
People say a chatbot ``understands language.''
Under the hood, the deployed object is a parameterized function that maps a prefix of tokens to a distribution over the next token.

\problem
Language is open-ended. Hand-written rules (parsers + templates) do not scale to web text, code, dialogue, and instructions.

\rootcause
We need one objective that (i)~uses abundant unlabeled text and (ii)~produces a generative system.

\solution
\textbf{Autoregressive language modeling.}
Model a sequence by the chain rule
\begin{equation}
P(x_1,\ldots,x_T)=\prod_{t=1}^{T} P(x_t \mid x_{<t}),
\end{equation}
and train by minimizing next-token cross-entropy.
At inference, sample or take \(\arg\max\) tokens left to right~\cite{brown2020language}.

\payoff
One scalable objective can absorb many downstream behaviors as ``continue this prompt usefully.''

\subsection{Why Transformers?}
\label{sec:transformer}

\setup
We need an architecture that estimates \(P(x_t\mid x_{<t})\) for long contexts and trains efficiently on accelerators.

\problem
RNNs process tokens mostly sequentially; long-range dependencies and GPU parallelism are awkward at web scale.

\rootcause
The architecture must let every position gather information from other positions in parallel, with enough capacity to store patterns.

\solution
The \textbf{Transformer}~\cite{vaswani2017attention}: self-attention (Query/Key/Value) plus position-wise feed-forward networks.
Decoder-only LLMs add \textbf{causal masking} so position \(t\) cannot attend to future tokens.

\payoff
Expressivity plus hardware fit made Transformers the default LLM backbone.

\subsection{Tokenization}
\label{sec:tok}

\setup
Neural nets consume integer token IDs, not raw Unicode strings.

\problem
Character-level sequences are too long (attention cost).
Word-level vocabularies explode; rare words, typos, and code break.

\rootcause
We need an open vocabulary without pathological sequence length.

\solution
\textbf{Subword tokenization}, commonly Byte-Pair Encoding (BPE) / byte-level BPE in GPT-style models; other families use WordPiece or Unigram/SentencePiece.

\payoff
Any string is encodable; frequent words stay compact; rare pieces split.

\subsection{Scaling laws}
\label{sec:scaling}

\setup
Pretraining spend must choose model size and token count under a FLOP budget.

\problem
``Just make it bigger'' can waste compute if the model is undertrained (or over-wide relative to data).

\rootcause
Loss depends jointly on parameters, data, and compute---not on parameters alone.

\solution
Empirical scaling laws~\cite{kaplan2020scaling,hoffmann2022chinchilla}.
Chinchilla emphasizes compute-optimal tokens-per-parameter allocations.

\payoff
Guides budgets.
\textbf{Caution:} fits are empirical; data quality shifts constants; downstream skills need not scale smoothly.

% =============================================================================
\section{Pretraining}
\label{sec:pretrain}

\subsection{Pretraining data}
\label{sec:data}

\setup
Pretraining is next-token prediction on massive corpora (web, books, code, \ldots).

\problem
Raw crawls contain spam, boilerplate, duplicates, and PII.
Duplicates waste capacity and raise contamination risk.

\rootcause
Gradients follow the empirical data distribution.

\solution
Filter, deduplicate, and \textbf{mix domains}.
Modern open recipes (e.g., Llama-family reports) describe these steps.
Synthetic rewritten data appears carefully; overuse risks collapse.

\payoff
Higher effective quality per token; stronger code/math when those domains are present.

\subsection{What pretraining does not guarantee}
\label{sec:pretrain-limits}

\setup
After pretraining you have a \textbf{base model}.

\problem
Base models complete text; they need not follow instructions, refuse safely, or prefer helpful answers among many fluent ones.

\rootcause
The objective was ``predict the next token on internet-like text,'' not ``be a helpful assistant.''

\solution
Hand the baton to post-training (Section~\ref{sec:post}).
Prompting cannot fully invent missing base competence; pretraining alone does not create chat alignment.

\subsection{Architecture and systems knobs}
\label{sec:arch-sys}

Explain each the same way (setup $\to$ problem $\to$ root cause $\to$ solution).
Compressed catalog:

\begin{table}[ht]
\centering
\small
\begin{tabularx}{\textwidth}{@{}lYY@{}}
\toprule
\textbf{Method} & \textbf{Problem targeted} & \textbf{Idea} \\
\midrule
RMSNorm / pre-norm & Deep-net instability & Normalize inside blocks \\
SwiGLU & Need MLP quality per param & Gated FFN \\
RoPE & Absolute PE weak for length & Rotary relative positions \\
GQA / MQA & KV bandwidth at decode & Share KV heads \\
MoE & Dense compute too expensive & Sparse experts per token \\
AdamW + schedules & Stable large-batch training & Adaptive opt.\ + warmup/decay \\
Mixed precision & Throughput / memory walls & BF16/FP16/FP8 matmuls \\
DP / TP / PP / FSDP / ZeRO & Model+optimizer do not fit one GPU & Shard state and compute \\
\bottomrule
\end{tabularx}
\caption{Pretraining architecture and systems methods (interview stubs).}
\label{tab:pre-knobs}
\end{table}

% =============================================================================
\section{Post-Training}
\label{sec:post}

\subsection{Supervised fine-tuning (SFT)}
\label{sec:sft}

\setup
Users want answers to instructions, not random web continuations.

\problem
A base model may continue like a webpage, ignore the ask, or emit the wrong format.

\rootcause
Pretraining distribution \(\neq\) assistant distribution.

\solution
Fine-tune on \((\mathrm{prompt},\mathrm{desired\ response})\) demonstrations with ordinary likelihood.
This is Step~1 of InstructGPT~\cite{ouyang2022instructgpt}.

\payoff
Teaches format and basic instruction following.
Still insufficient when many answers are fluent but differently preferred.

\subsection{Classic RLHF --- do not start with ``PPO''}
\label{sec:rlhf}

\setup
After SFT, multiple fluent answers exist; humans still have preferences.

\problem
SFT clones demonstrations; it does not directly optimize ``prefer \(A\) over \(B\).''
Writing one gold string per open-ended prompt is incomplete.

\rootcause
We need a ranking signal, then policy optimization that does not reward-hack into gibberish.

\solution
InstructGPT's three steps~\cite{ouyang2022instructgpt,openai2022instruct}:
\begin{enumerate}[leftmargin=*]
  \item SFT on demonstrations;
  \item train a \textbf{reward model} on human comparisons of model outputs;
  \item optimize the policy with \textbf{PPO} to raise RM reward, with a \textbf{KL penalty} to a reference policy.
\end{enumerate}

\payoff
First widely deployed instruction-aligned GPT-style assistants.
Costs: preference data, RM misspecification, reward hacking, multi-model complexity.

\subsection{RLAIF / Constitutional AI}
\label{sec:rlaif}

\setup
Human labels do not scale forever.

\problem
Safety/helpfulness coverage is expensive and thin on edges.

\rootcause
Human feedback is the bottleneck.

\solution
AI feedback (RLAIF) and principle-guided critique/revision (Constitutional AI)~\cite{bai2022constitutional}.

\payoff
Scales labeling; inherits teacher-model biases.

\subsection{DPO --- problem first}
\label{sec:dpo}

\setup
Full RLHF needs sampling, an RM, PPO, a value model, and KL control.

\problem
That stack is costly and operationally brittle for many labs.

\rootcause
If preferences already define a Bradley--Terry ranking, the KL-constrained optimal policy can be expressed so an explicit RM+PPO loop is not mandatory.

\solution
\textbf{Direct Preference Optimization (DPO)}~\cite{rafailov2023dpo}: train the policy on preference pairs with an objective that implicitly encodes the reward.
Related methods (IPO, KTO, ORPO, SimPO, \ldots) should be introduced only after naming which RLHF pain point they target.

\payoff
Simpler alignment; still depends on preference quality.

\subsection{Reasoning post-training}
\label{sec:reasoning}

\setup
Chat alignment \(\neq\) contest math/coding reliability.

\problem
Outcome-only rewards are sparse; soft learned rewards can be hacked; PPO critics are heavy.

\rootcause
Credit assignment and reward design for long traces are hard.

\solution
\begin{itemize}[leftmargin=*]
  \item ORM: score final answers when verifiers exist;
  \item PRM: score intermediate steps (process supervision);
  \item RL with verifiable rewards (tests/checkers);
  \item \textbf{GRPO}~\cite{shao2024deepseekmath}: group of samples per prompt; group-normalized advantages; no value network;
  \item DeepSeek-R1 multi-stage recipe~\cite{deepseekr12025}: R1-Zero uses rule-based RL; full R1 adds cold-start SFT and further stages for readability/helpfulness.
\end{itemize}

\payoff
Can raise reasoning benchmarks when rewards are grounded.
Do not invent exact leaderboard numbers; copy from a specific paper revision.

\subsection{Parameter-efficient fine-tuning}
\label{sec:peft}

\setup
Full fine-tunes are memory-heavy; many tasks may need adapters.

\problem
A full weight copy per task is impractical.

\rootcause
Task updates are often low-dimensional relative to \(W\).

\solution
LoRA, QLoRA, adapters, prefix/prompt tuning, model merging.

\payoff
Fine-tune on smaller GPUs; serve multiple adapters.

% =============================================================================
\section{Inference}
\label{sec:inference}

\subsection{Prefill vs decode}
\label{sec:prefill-decode}

\setup
A user sends a prompt; the system returns tokens.

\problem
Calling everything ``inference'' hides the structure serving engineers use~\cite{vizuara2026paged}:
\begin{enumerate}[leftmargin=*]
  \item \textbf{Prefill:} process prompt tokens (largely parallel) and emit the first output token;
  \item \textbf{Decode:} emit one new token at a time.
\end{enumerate}

\rootcause
Autoregressive factorization forces sequential decode even when prefill is parallel.

\payoff
KV cache, continuous batching, and PagedAttention are all explained relative to this split.

\subsection{KV cache}
\label{sec:kv}

\setup
Each decode step needs keys/values for all previous tokens.

\problem
Recomputing \(K/V\) for the entire prefix every step repeats matmuls; latency grows badly with length.

\rootcause
Past tokens' \(K/V\) do not change; only the new token adds new \(K/V\).

\solution
\textbf{KV cache:} store past keys/values; append each new token's \(K/V\).

\payoff
Less compute, more memory---and KV memory becomes the serving bottleneck that motivates PagedAttention.

\subsection{PagedAttention (full interview-style walkthrough)}
\label{sec:paged}

This subsection follows the layered answer structure taught in~\cite{vizuara2026paged}, with mechanisms validated against the vLLM writeup and paper~\cite{kwon2023vllm,vllmblog2023}.

\subsubsection*{Layer 1 --- Basics: GPU memory during inference}
GPU memory during inference holds roughly:
(1)~model weights,
(2)~activations,
(3)~the KV cache, which grows with generated tokens and often dominates serving.

\subsubsection*{Layer 2 --- Problem: contiguous KV reservation wastes memory}
The KV cache is large and \emph{dynamic} (length unknown ahead of time).
The vLLM authors report that existing systems can waste a large fraction of KV memory via fragmentation and over-reservation (characterized as about 60\%--80\% waste in the regimes discussed in their blog)~\cite{vllmblog2023}.

Teaching example with multiple users~\cite{vizuara2026paged}:
reserve a contiguous slab per request up to a projected \(\texttt{max\_tokens}\).
Then:
\begin{itemize}[leftmargin=*]
  \item reserved space sits unused while a request is still decoding, blocking other users;
  \item when a short request finishes early, free holes may be too small or wrongly placed for a waiting longer request.
\end{itemize}
Effective batch size collapses; throughput falls.

\subsubsection*{Layer 3 --- Root cause}
KV for a request is allocated as contiguous physical memory, despite highly variable lengths---analogous to forcing contiguous physical RAM for an entire process address space.

\subsubsection*{Layer 4 --- Solution: PagedAttention in vLLM}
Inspired by OS virtual memory and paging~\cite{kwon2023vllm,vllmblog2023}:
\begin{enumerate}[leftmargin=*]
  \item Partition KV memory into fixed-size \textbf{blocks} (pages). Block size is a system parameter; discussions of vLLM often use 16 tokens/block as a default example.
  \item Each block stores \(K/V\) for that many tokens.
  \item Maintain a \textbf{free block list} and a per-request \textbf{block table} mapping logical spans to physical blocks.
  \item Allocate blocks on demand during decode; a request's blocks need not be contiguous.
  \item The attention kernel gathers \(K/V\) by following the block table (\textbf{PagedAttention}).
\end{enumerate}
The same abstraction enables sharing blocks across sequences (e.g., parallel samples from one prompt) via reference counting / copy-on-write~\cite{vllmblog2023}.

\subsubsection*{Layer 5 --- Payoff}
Waste is largely limited to the last partial block (blog characterization: under \(\sim\)4\% in their analysis)~\cite{vllmblog2023}, enabling larger batches and higher throughput.
Reported speedups are workload-dependent; do not universalize a single multiplier.

\subsection{Decoding, systems, and behavioral inference methods}
\label{sec:decode-behave}

\begin{table}[ht]
\centering
\small
\begin{tabularx}{\textwidth}{@{}lYY@{}}
\toprule
\textbf{Method} & \textbf{Problem} & \textbf{Root idea} \\
\midrule
Temperature / top-\(k\) / top-\(p\)~\cite{holtzman2019curious} & Greedy dullness vs long-tail nonsense & Decision policy on \(P(x_t\mid x_{<t})\) \\
Continuous batching & Static batches idle as sequences finish unevenly & Dynamically admit/finish requests \\
Speculative decoding~\cite{leviathan2023speculative} & Large-model serial decode latency & Draft proposes; large verifies \\
Quantization & VRAM / bandwidth walls & Lower-bit weights/KV \\
Cascade / routing~\cite{ong2025routellm,ding2024hybrid} & Always-large is too costly & Easy queries $\to$ smaller models \\
CoT~\cite{wei2022cot} / self-consistency~\cite{wang2022selfconsistency} & Multi-step one-shot failures & Extra reasoning tokens / vote \\
RAG~\cite{lewis2020rag} / tools & Stale parametric memory; weak exact ops & Retrieve or call external systems \\
Test-time compute & Residual hard-tail errors & Spend more FLOPs at serve time \\
\bottomrule
\end{tabularx}
\caption{Inference methods, each still explained via problem-first narration.}
\label{tab:infer-methods}
\end{table}

% =============================================================================
\section{Evaluation}
\label{sec:eval}

\subsection{Why evaluation is its own stage}
\label{sec:eval-why}

\setup
You changed data, loss, or decoding; a dashboard number moved.

\problem
A leaderboard can rise while users get worse or less safe answers.

\rootcause
Benchmarks operationalize \emph{narrow} skills; proxies can be gamed; contamination inflates scores; LLM-as-judge has biases~\cite{zheng2023judging}.

\solution
Portfolio evaluation: capability suites (e.g., MMLU~\cite{hendrycks2021mmlu}, GSM8K~\cite{cobbe2021gsm8k}, HumanEval~\cite{chen2021humaneval}), instruction rubrics, human preference (Arena Elo), safety, plus cost/latency.
Report threats: contamination, prompt sensitivity, judge bias, shift.

\payoff
Know whether the change was real.

\subsection{Routing metrics}
\label{sec:routing-eval}

\setup
Choose weak/cheap vs strong/expensive per query~\cite{ong2025routellm,ding2024hybrid}.

\problem
Accuracy alone hides the point of routing.

\rootcause
The objective is Pareto: quality vs cost.

\solution
\begin{equation}
\mathrm{PGR}(R)=\frac{r(R)-r(M_{\mathrm{weak}})}{r(M_{\mathrm{strong}})-r(M_{\mathrm{weak}})},
\end{equation}
plus cost \(c(R)\) and CPT(\(x\%\)) (minimum strong-call fraction to hit target PGR).

% =============================================================================
\section{End-to-End Causal Chain}
\label{sec:chain}

\begin{lstlisting}[basicstyle=\ttfamily\small,frame=single]
Text in the world
  -> filter/dedup/mix + tokenize
  -> pretrain next-token LM
  -> base model
  -> SFT
  -> preferences (RLHF or DPO/...) / reasoning RL
  -> aligned / reasoning model
  -> serve: prefill/decode, KV, batching, PagedAttention, sampling
  -> optionally: RAG, tools, test-time compute, routing
  -> evaluate capability + preference + safety + cost
  -> iterate
\end{lstlisting}

\begin{table}[ht]
\centering
\small
\begin{tabularx}{\textwidth}{@{}lY@{}}
\toprule
\textbf{Symptom} & \textbf{Likely stage} \\
\midrule
Cannot do basics even with good prompts & Pretraining data/capacity \\
Knows facts, ignores instructions & SFT \\
Fluent but rude/unsafe/sycophantic & Preference / safety post-training \\
Chat OK, hard math fails & Reasoning post-training or test-time compute \\
Correct but too slow/expensive & Inference systems / routing \\
Leaderboard up, users unhappy & Evaluation mismatch \\
\bottomrule
\end{tabularx}
\caption{Symptom to lifecycle stage.}
\label{tab:symptoms}
\end{table}

% =============================================================================
\section{Interview Cheat Sheet}
\label{sec:cheat}

When asked ``What is \(X\)?'':
\begin{enumerate}[leftmargin=*]
  \item Do \textbf{not} start with the definition of \(X\).
  \item Explain system context (training or serving).
  \item Show the naive approach and where it breaks.
  \item Name the root design mistake.
  \item Introduce \(X\) and walk the mechanism.
  \item State the payoff and the new tradeoff.
  \item Cite one credible source.
\end{enumerate}
That is the PagedAttention-video pattern~\cite{vizuara2026paged}, applied to every concept above.

% =============================================================================
\section{Explicit Non-Claims}
\label{sec:nonclaims}

\begin{itemize}[leftmargin=*]
  \item No universal ``best'' alignment algorithm.
  \item No fabricated benchmarks or universalized speedups.
  \item Preference alignment does not ``solve'' safety.
  \item Agents add orchestration/tools/memory \emph{on top of} this stack.
\end{itemize}

% =============================================================================
\begin{thebibliography}{99}

\bibitem{vizuara2026paged}
Vizuara.
\newblock LLM Interview Series \#5: What Is PagedAttention?
\newblock YouTube, \url{https://www.youtube.com/watch?v=-AB6m0Spo6c}, 2026.
\newblock Used here as a \emph{narration-style} reference, not as a primary technical source.

\bibitem{vaswani2017attention}
Ashish Vaswani et~al.
\newblock Attention is all you need.
\newblock In \emph{NeurIPS}, 2017.

\bibitem{brown2020language}
Tom~B. Brown et~al.
\newblock Language models are few-shot learners.
\newblock In \emph{NeurIPS}, 2020.

\bibitem{kaplan2020scaling}
Jared Kaplan et~al.
\newblock Scaling laws for neural language models.
\newblock arXiv:2001.08361, 2020.

\bibitem{hoffmann2022chinchilla}
Jordan Hoffmann et~al.
\newblock Training compute-optimal large language models.
\newblock In \emph{NeurIPS}, 2022.

\bibitem{ouyang2022instructgpt}
Long Ouyang et~al.
\newblock Training language models to follow instructions with human feedback.
\newblock In \emph{NeurIPS}, 2022.

\bibitem{openai2022instruct}
OpenAI.
\newblock Aligning language models to follow instructions.
\newblock \url{https://openai.com/index/instruction-following/}, 2022.

\bibitem{bai2022constitutional}
Yuntao Bai et~al.
\newblock Constitutional AI: Harmlessness from AI feedback.
\newblock arXiv:2212.08073, 2022.

\bibitem{rafailov2023dpo}
Rafael Rafailov et~al.
\newblock Direct preference optimization: Your language model is secretly a reward model.
\newblock In \emph{NeurIPS}, 2023.

\bibitem{shao2024deepseekmath}
Zhihong Shao et~al.
\newblock DeepSeekMath: Pushing the limits of mathematical reasoning in open language models.
\newblock arXiv:2402.03300, 2024.

\bibitem{deepseekr12025}
DeepSeek-AI.
\newblock DeepSeek-R1: Incentivizing reasoning capability in LLMs via reinforcement learning.
\newblock 2025.

\bibitem{kwon2023vllm}
Woosuk Kwon et~al.
\newblock Efficient memory management for large language model serving with PagedAttention.
\newblock In \emph{SOSP}, 2023.

\bibitem{vllmblog2023}
Woosuk Kwon and Zhuohan Li.
\newblock vLLM: Easy, fast, and cheap LLM serving with PagedAttention.
\newblock \url{https://vllm.ai/blog/2023-06-20-vllm}, 2023.

\bibitem{holtzman2019curious}
Ari Holtzman et~al.
\newblock The curious case of neural text degeneration.
\newblock In \emph{ICLR}, 2020.

\bibitem{leviathan2023speculative}
Yaniv Leviathan, Matan Kalman, and Yossi Matias.
\newblock Fast inference from transformers via speculative decoding.
\newblock In \emph{ICML}, 2023.

\bibitem{lewis2020rag}
Patrick Lewis et~al.
\newblock Retrieval-augmented generation for knowledge-intensive NLP tasks.
\newblock In \emph{NeurIPS}, 2020.

\bibitem{wei2022cot}
Jason Wei et~al.
\newblock Chain-of-thought prompting elicits reasoning in large language models.
\newblock In \emph{NeurIPS}, 2022.

\bibitem{wang2022selfconsistency}
Xuezhi Wang et~al.
\newblock Self-consistency improves chain of thought reasoning in language models.
\newblock In \emph{ICLR}, 2023.

\bibitem{hendrycks2021mmlu}
Dan Hendrycks et~al.
\newblock Measuring massive multitask language understanding.
\newblock In \emph{ICLR}, 2021.

\bibitem{cobbe2021gsm8k}
Karl Cobbe et~al.
\newblock Training verifiers to solve math word problems.
\newblock arXiv:2110.14168, 2021.

\bibitem{chen2021humaneval}
Mark Chen et~al.
\newblock Evaluating large language models trained on code.
\newblock arXiv:2107.03374, 2021.

\bibitem{zheng2023judging}
Lianmin Zheng et~al.
\newblock Judging LLM-as-a-judge with MT-Bench and Chatbot Arena.
\newblock In \emph{NeurIPS}, 2023.

\bibitem{ong2025routellm}
Isaac Ong et~al.
\newblock RouteLLM: Learning to route LLMs with preference data.
\newblock In \emph{ICLR}, 2025.

\bibitem{ding2024hybrid}
Dujian Ding et~al.
\newblock Hybrid LLM: Cost-efficient and quality-aware query routing.
\newblock arXiv:2404.14618, 2024.

\end{thebibliography}

\end{document}
